\section*{Глава 3. Разработка веб-приложения}
\addcontentsline{toc}{section}{Глава 3. Разработка веб-приложения}
\stepcounter{section}

\subsection{Разработка базовой структуры приложения и настройка окружения}

В соответствии с задачей 3.1 было разработано веб-приложение для управления личным гардеробом и формирования пользовательских образов. Для серверной части был выбран фреймворк \textbf{Flask}, так как он позволяет быстро реализовать маршрутизацию, обработку форм, работу с шаблонами и подключение базы данных.

Структура проекта включает основной файл приложения \texttt{app.py}, папку HTML-шаблонов \texttt{templates}, статические файлы \texttt{static}, папку для загружаемых изображений \texttt{static/uploads}, а также файл зависимостей \texttt{requirements.txt}. В качестве шаблонизатора используется \textbf{Jinja2}, что позволяет формировать страницы на стороне сервера и передавать в них данные из базы.

Для хранения данных используется реляционная база данных. В процессе разработки применяется SQLite, так как она не требует отдельной установки сервера БД. При этом структура моделей реализована через \textbf{SQLAlchemy}, что позволяет при необходимости подключить PostgreSQL через переменную окружения \texttt{DATABASE\_URL}.

Базовая инициализация приложения представлена в листинге~\ref{lst:flask_init}.

\begin{lstlisting}[caption={Инициализация Flask-приложения}, label={lst:flask_init}]
app = Flask(__name__)
app.config["SQLALCHEMY_DATABASE_URI"] = database_url
db = SQLAlchemy(app)
\end{lstlisting}

Таким образом, на данном этапе была подготовлена основная структура веб-приложения, подключены зависимости, настроена работа с базой данных и организовано хранение статических файлов.

\subsection{Реализация аутентификации, авторизации и ролей пользователей}

В соответствии с задачей 3.2 в приложении реализована система регистрации, входа и выхода из учетной записи. Пользователь может создать аккаунт, после чего его данные сохраняются в таблице \texttt{users}. Пароль не хранится в открытом виде: перед сохранением он преобразуется в хеш.

В системе предусмотрены роли пользователей: обычный пользователь \texttt{user} и администратор \texttt{admin}. Обычный пользователь получает доступ к личному гардеробу и созданным образам. Администратор имеет доступ к панели управления, где может просматривать пользователей, изменять роли, управлять категориями и выполнять модерацию данных.

Проверка доступа к защищенным страницам выполняется с помощью декораторов. Пример проверки роли администратора представлен в листинге~\ref{lst:role_required}.

\begin{lstlisting}[caption={Проверка роли пользователя}, label={lst:role_required}]
def role_required(role):
    def decorator(view):
        @wraps(view)
        def wrapped_view(*args, **kwargs):
            user = current_user()
            if user is None:
                flash("Для доступа к странице необходимо войти в систему.", "warning")
                return redirect(url_for("login", next=request.path))
            if user.role != role:
                abort(403)
            return view(*args, **kwargs)

        return wrapped_view

    return decorator
\end{lstlisting}

Таким образом, в приложении реализовано разграничение прав доступа, что соответствует требованиям к нескольким видам пользователей.

\subsection{Реализация CRUD-интерфейса управления гардеробом}

В соответствии с задачей 3.3 был реализован CRUD-интерфейс для управления элементами гардероба. Пользователь может добавлять новые предметы одежды, просматривать список вещей, редактировать информацию о предмете и удалять его из системы.

Для каждого элемента одежды сохраняются следующие данные: название, категория, цвет, сезон и путь к изображению. Категории хранятся в отдельной таблице \texttt{categories}, что позволяет избежать дублирования названий категорий в базе данных.

На странице гардероба реализованы фильтры по названию, категории, сезону и цвету. Для выбора цвета используется цветовой круг, что делает интерфейс более наглядным и удобным по сравнению с ручным вводом цвета.

Добавление нового предмета одежды выполняется через HTML-форму. После отправки формы данные проверяются, изображение сохраняется в папку \texttt{static/uploads}, а запись о предмете добавляется в таблицу \texttt{clothes}.

\begin{lstlisting}[caption={Создание элемента гардероба}, label={lst:create_clothing}]
item = ClothingItem(
    user_id=current_user().id,
    category_id=category.id,
    name=name,
    color=request.form.get("color"),
    season=normalize_season(request.form.get("season")),
    image_filename=image_filename,
)

db.session.add(item)
db.session.commit()
\end{lstlisting}

При удалении предмета одежды приложение также удаляет его изображение и обновляет связанные образы, чтобы в них не оставались ссылки на удалённый элемент.

\subsection{Реализация формирования и управления пользовательскими образами}

В соответствии с задачей 3.4 была реализована возможность создания пользовательских образов на основе предметов гардероба. Пользователь может выбрать несколько вещей, указать название образа и сохранить комплект. Связь между образом и элементами одежды реализована через промежуточную таблицу \texttt{outfit\_items}, что соответствует связи «многие ко многим».

Также реализована автоматическая генерация образа. Пользователь выбирает сезон и любимый цвет, после чего приложение подбирает подходящие вещи из гардероба. Алгоритм учитывает категорию предмета, сезонность и выбранный цвет. Если в гардеробе есть вещь любимого цвета, она обязательно включается в сгенерированный образ.

Образ формируется из нескольких зон: верх, низ, верхняя одежда, обувь, аксессуары и головные уборы. Такая структура позволяет собрать комплект, близкий к реальному образу пользователя.

\begin{lstlisting}[caption={Формирование связи образа с вещами}, label={lst:outfit_items}]
outfit = Outfit(
    user_id=user.id,
    name=form_values["name"],
)

db.session.add(outfit)
db.session.flush()
outfit.items = selected_items
db.session.commit()
\end{lstlisting}

Пользователь может просматривать сохранённые образы, редактировать их состав, удалять ненужные комплекты и скачивать итоговый коллаж в формате PNG.

\subsection{Реализация загрузки, хранения и выгрузки изображений}

В соответствии с задачей 3.5 была реализована загрузка, хранение и выгрузка изображений. При добавлении предмета одежды пользователь загружает файл изображения. Приложение проверяет расширение файла и корректность изображения, после чего сохраняет его в папку \texttt{static/uploads}. В базе данных хранится не само изображение, а путь к нему.

Для выгрузки отдельного изображения используется отдельный маршрут, который отдаёт файл пользователю через функцию \texttt{send_file}. Кроме того, для образов реализована выгрузка коллажа. Коллаж формируется автоматически на белом фоне с помощью библиотеки \textbf{Pillow}. Изображения одежды размещаются по заранее заданным зонам в зависимости от категории предмета.

\begin{lstlisting}[caption={Создание PNG-коллажа образа}, label={lst:collage}]
canvas = Image.new("RGBA", (1200, 1200), "white")

with Image.open(image_path) as image:
    paste = image.convert("RGBA")
    paste.thumbnail((max_width, max_height))
    canvas.alpha_composite(paste, (paste_x, paste_y))

buffer = BytesIO()
canvas.convert("RGB").save(buffer, format="PNG")
buffer.seek(0)
\end{lstlisting}

После формирования изображение не сохраняется на диск, а передаётся пользователю напрямую из памяти. Это позволяет каждый раз создавать актуальный коллаж на основе текущего состава образа.

\begin{lstlisting}[caption={Выгрузка коллажа пользователю}, label={lst:export_collage}]
return send_file(
    buffer,
    mimetype="image/png",
    as_attachment=True,
    download_name=f"outfit-{outfit.id}.png"
)
\end{lstlisting}

Таким образом, в третьем разделе были реализованы основные функции веб-приложения: структура проекта, регистрация и авторизация пользователей, разграничение прав доступа, управление гардеробом, создание пользовательских образов, загрузка изображений и выгрузка готового коллажа.
